\documentclass{ieeeaccess}
% IEEE Access manuscript for Study 001.
% Compile with pdflatex (twice) from this directory; the class, fonts, and
% logos live in ./ieeeaccess_support/. Adjust \graphicspath if you move files.
% All figure files are PLACEHOLDERS (see the % TODO: replace ... captions).
% Author/affiliation blocks are placeholders — fill before submission.
% Status (2026-08-06): GPU-free AE revision — real workflow figures and tables
% wired in (design/dag/pipeline/lineage/G3-trace/D3), Related Work expanded with
% verified 2026 references, appendices restructured (identification argument,
% D3 battery, Woodbury/REML, cost plan + roster). Remaining TODOs: history dates,
% DOI, funding statement, biographies, and the pending-eval figures.
\usepackage{cite}
\usepackage{amsmath,amssymb,amsfonts}
\usepackage{graphicx}
\usepackage{textcomp}
\usepackage{url}

\graphicspath{{./ieeeaccess_support/}}

\usepackage{bm}
\makeatletter
\AtBeginDocument{\DeclareMathVersion{bold}
\SetSymbolFont{operators}{bold}{T1}{times}{b}{n}
\SetSymbolFont{NewLetters}{bold}{T1}{times}{b}{it}
\SetMathAlphabet{\mathrm}{bold}{T1}{times}{b}{n}
\SetMathAlphabet{\mathit}{bold}{T1}{times}{b}{it}
\SetMathAlphabet{\mathbf}{bold}{T1}{times}{b}{n}
\SetMathAlphabet{\mathtt}{bold}{OT1}{pcr}{b}{n}
\SetSymbolFont{symbols}{bold}{OMS}{cmsy}{b}{n}
\renewcommand\boldmath{\@nomath\boldmath\mathversion{bold}}}
\makeatother

\def\BibTeX{{\rm B\kern-.05em{\sc i\kern-.025em b}\kern-.08em
    T\kern-.1667em\lower.7ex\hbox{E}\kern-.125emX}}

% The IEEE Access class re-sets \ttdefault to Courier (pcr) at \begin{document};
% its T1 metric (pcrr8t) is not shipped with BasicTeX, so re-override it there.
\AtBeginDocument{\renewcommand{\ttdefault}{cmtt}}

\begin{document}
\history{Date of publication xxxx 00, 0000, date of current version xxxx 00, 0000. % TODO: set dates
}
\doi{10.1109/ACCESS.2024.0000000}% TODO: assigned at submission

\title{Lineage or Era? An Identifiability-Gated Variance Decomposition of Correlated Errors in Public Language Models}

\author{\uppercase{S. Kanaga Suba Raja}\authorrefmark{1},
\uppercase{Sudharsan S}\authorrefmark{1}}% TODO: confirm author order and any IEEE membership designations

\address[1]{Department of Computer Science and Engineering, SRM Institute of
Science and Technology, Tiruchirappalli, Tamil Nadu 621105, India (e-mail:
kanagass@srmist.edu.in; ss0856@srmist.edu.in)}% TODO: confirm postal codes and secondary emails

\tfootnote{This work was supported in part by ... .% TODO: funding/support statement
}

\markboth
{Raja and Sudharsan: Identifiability-Gated Variance Decomposition of Correlated Errors in Public Language Models}
{Raja and Sudharsan: Identifiability-Gated Variance Decomposition of Correlated Errors in Public Language Models}

\corresp{Corresponding author: S. Kanaga Suba Raja (e-mail: kanagass@srmist.edu.in).}

\begin{abstract}
Public language models make correlated errors: when one open-weight model fails a question, its contemporaries and descendants often fail it as well. Whether this shared failure structure is driven by lineage---a model replicating the blind spots of the models and data it descends from---or by era---models released together acquiring the same errors through shared training data and technique---is untested, because the two are confounded: descendants are released after their ancestors, making release year both a mediator and a confounder of the lineage path. We introduce an identifiability-gated variance-decomposition instrument, built on the documented release record rather than reconstructed family trees, that partitions correlated error into lineage, era, and model-specific components on the connected subset of the open-model population where the partition is identifiable. The design is gated in two stages. First, a simulation study validates that a crossed restricted-maximum-likelihood estimator recovers known ground truth under realistic occupancy and fails detectably when family and era are nested (aliased). Second, a pre-analysis study-population design procedure---which never observes trait values---selects the smallest population, 22 of 47 models, that remains structurally identifiable and clears a simulation-based recoverability bar. The empirical partition itself is not claimed here: it is the object of a pre-registered measurement pass that this paper scopes, gates, and cost-bounds, and whose lineage-versus-era shares will be reported against a pre-registered decision rule for whether cross-family diversification is a credible remedy for correlated failure.
\end{abstract}

\begin{keywords}
Correlated errors, error correlation, foundation models, identifiability, mixed-effects models, restricted maximum likelihood, variance components, variance decomposition.
\end{keywords}

\titlepgskip=-21pt

\maketitle

\section{Introduction}
\label{sec:introduction}
\PARstart{T}{he} growing reliance on a small number of public open-weight language
models raises a statistical problem that is older than the models themselves:
when a group of models is given a hard question, they often fail together---missing
the same fact, making the same mistake, reaching the same wrong answer. Two
explanations compete for this pattern of shared mistakes. The first is
\emph{lineage}: a model trained on another model's outputs, or on data generated by
it, tends to replicate that model's blind spots. The second is \emph{era}: models
released at the same time are trained on largely the same scrapes of the web, the
same benchmark-cleansed corpora, and the same dominant techniques, so they may pick
up the same errors independently of any shared ancestry \cite{kim2025}.

Whether shared failure is mostly inherited or mostly contemporaneous is not an
academic question. If shared mistakes trace mostly to lineage, then
diversification across model families is a credible mitigation: deploying models
with independent ancestries reduces the probability that the ensemble fails a
question the same way. If they trace mostly to era, then no amount of new ancestry
helps, because every model of a given year carries the same errors. The choice
between these two conclusions has direct operational consequences for how
practitioners assemble model portfolios.

The difficulty is that the two explanations are not separable by inspection.
Descendants are always released after their ancestors, so release year lies on
the lineage-to-error path: a child model inherits both the errors of its parent
and the training environment of its own release date. Release year is therefore a
\emph{mediator} of the lineage path and a \emph{potential confounder} of the
observational family--error association (the formal argument is in
Appendix~\ref{app:identification}). Any single-number answer that tries to assign
shared error ``to lineage'' or ``to era'' conflates the two roles of release year.

We address the problem with a variance-decomposition instrument rather than with a
causal attribution. Treating the open-model population as a designed experiment, we
partition the observed pattern of shared errors into a lineage component
$\sigma^{2}_{L}$, an era component $\sigma^{2}_{E}$, and a model-specific remainder
$\sigma^{2}_{U}$, estimated with a crossed random-effects (restricted-maximum-likelihood)
model. Because the partition is defined only where lineage and era actually vary
together---the \emph{connected subset} of the population where the design is
crossed rather than nested---the analysis is explicitly scoped to that subset, and
an identifiability gate decides whether the decomposition is runnable at all
\cite{harville1977,patterson1971}.

The central methodological claim of this paper is the ordering of decisions.
Identifiability is verified before any trait is measured; the estimator is
validated in simulation before any real data are analyzed; and the study population
is selected by a pre-analysis design procedure that never observes trait values,
before any empirical inference. Section \ref{sec:questions} states the research
questions; Section \ref{sec:related} situates the work; Section
\ref{sec:model} defines the estimands and identifiability conditions; Section
\ref{sec:methods} describes the three stages of the protocol (population audit,
simulation validation, pre-analysis population design); Section \ref{sec:results}
reports the validated results and marks the empirical outputs that are pending the
measurement campaign; Sections \ref{sec:discussion} and \ref{sec:threats} discuss
implications and threats to validity.

\section{Problem Statement and Research Questions}
\label{sec:questions}

\subsection{Main question}
The main research question (RQ0) asks whether the observed error-correlation
structure across public open-weight language models can be separated into a lineage
component and an era component; that is, whether the partition
$\sigma^{2}_{L}+\sigma^{2}_{E}+\sigma^{2}_{U}$ is estimable and non-degenerate on
the connected subset. A positive answer requires both: (a) the estimator survives
simulation validation under the realistic population occupancy, and (b) on real
data, the lineage and era components are each estimated with non-degenerate
intervals.

\subsection{Sub-questions}
Table~\ref{tab:rqs} lists the sub-questions, their operating population, and their
refutation conditions. RQ3 is the primary observational estimand (lineage
conditional on era); RQ4 is the mechanistic estimand (lineage with era held fixed),
available only on the small co-released and staggered-fine-tune subsets and reported
separately. RQ2 is the stop-rule question: a mis-specified generative model must
fail \emph{detectably} rather than silently, or the program stops.

\begin{table}[t]
\caption{\textbf{Research questions and refutation conditions.}}
\label{tab:rqs}
\setlength{\tabcolsep}{4pt}
\begin{tabular}{|p{22pt}|p{130pt}|p{70pt}|p{78pt}|}
\hline
RQ & Question & Operates on & Refutation \\
\hline
RQ1 & Does the crossed random-effects estimator recover known ground truth under a balanced-crossed (D1) and a realistic-occupancy (D2) regime? & Simulated data & Bias or collapse under D2 \\
RQ2 & Under a nested (mis-specified) design, does the estimator fail detectably? & Simulated data & Silent bias \\
RQ3 & What share of error covariance is attributable to lineage conditional on era? & Connected subset & $\sigma^{2}_{L}\approx0$ \\
RQ4 & Holding era exactly fixed, what is the structural contribution of lineage? & Co-released cohorts + fine-tune chains & Reported separately; never merged with RQ3 \\
RQ5 & After lineage adjustment, does the era component converge across release quarters? & Connected subset & Flat trend \\
RQ6 & Is cross-family diversification a credible mitigation for correlated failure? & Decision layer & Era dominates \\
\hline
\end{tabular}
\end{table}

\subsection{The two-estimand rule}
RQ3 and RQ4 answer different questions and are never merged. Release year is a
mediator of the lineage-to-error path and a potential confounder of the
observational family--error association, so the observational primary
adjusts for era grouping but cannot hold era fixed; the mechanistic secondary is the
only estimand that isolates lineage with era held exactly fixed, and it exists only
where co-released family cohorts and verified cross-generation fine-tune chains
provide the contrast. We report the two estimands as $\theta_{P}$ (observational)
and $\theta_{M}$ (mechanistic), separately, throughout.

\section{Related Work}
\label{sec:related}
Table~\ref{tab:diff} positions this work against its closest neighbors; this
section expands on the three nearest ones.

\paragraph{Agreement and co-failure.}
Kim \emph{et al.} measure cross-sectional agreement between language models---how
often two models agree when both are wrong---across more than 350 open and hosted
models on two leaderboards, finding roughly 60\% agreement among co-erring models
and agreement that persists across distinct architectures and providers
\cite{kim2025}. They are explicitly cross-sectional and leave causal and temporal
attribution open; they measure a correlation, not a variance partition. On the
failure side, Chen shows that pairwise error correlation substantially
\emph{underestimates} the co-failure ceiling of an ensemble---the probability that
every member fails a question---across 67 frontier models and three combination
schemes (routing, voting, mixture-of-agents), by roughly a factor of two
\cite{chen2026}. The operative risk is therefore the joint all-wrong rate, whose
governing inputs are exactly the lineage and era shares we estimate: our partition
quantifies the two components whose sum bounds that ceiling on the connected
population, so the two results are complementary rather than competing.

\paragraph{Variance decomposition of evaluation pipelines.}
The closest statistical neighbor is the total-evaluation-error (TEE) line
\cite{messing2026,brennan2001}. Messing decomposes the uncertainty of LLM
evaluation pipelines---judge-model choice, temperature, and prompt phrasing---into
crossed random effects via generalizability theory, and shows that naive standard
errors are 40--60\% smaller than the TEE-corrected ones and that naive confidence
intervals under-cover as $n$ grows \cite{messing2026}. The estimator class is the
same (crossed random effects, REML), but the grouping factors are properties of
the \emph{measurement pipeline}, not of the models; identifiability of a
model-population partition is never the question. The two objects are composable:
our trait-level decomposition sits on top of a pipeline whose measurement variance
TEE-style designs quantify. Classical generalizability theory applies to facets of
a measurement procedure rather than to a population of models whose
identifiability must be established first \cite{brennan2001}.

\paragraph{Lineage as structure.}
Phylogeny-reconstruction work treats lineage as a tree to be recovered from output
similarity: PhyloLM builds phylogenetic trees across open and closed models and
shows that tree distance predicts benchmark performance \cite{phylolm2024}. We
instead take the documented release record as given and treat lineage as a
variance component of a designed population. Tracing the Roots reconstructs the
evolution of \emph{datasets} in post-training---83 seed datasets giving rise to
430 datasets and 971 inheritance edges---and shows that benchmark contamination
propagates along lineage paths \cite{li2026roots}. Dataset ancestry is a channel
into our era component (shared training data is a shared environment), and its
finding that contamination travels along lineage is exactly the mechanism that
motivates the era channel, not a competing estimator of model error covariance.

\paragraph{Instrument-level bias and monoculture.}
Preference leakage shows that LLM-as-judge ratings are biased toward generators
that are the same model as, inherit from, or share a family with the judge
\cite{li2026pref}. This is a downstream \emph{measurement-instrument} consequence
of correlated error rather than a trait-level decomposition; it strengthens the
case for measuring error structure directly instead of through judge-mediated
leaderboards. The monoculture literature documents the systemic risk of deploying
many similar models but treats correlated error as a given input to a welfare
argument rather than as an object to be decomposed \cite{monoculture2024,jo2026};
Jo \emph{et al.} show that monoculture estimates are null-model-dependent and
argue that an estimand must be defined independently of a baseline, which is why
we model shared random effects directly rather than relying on an independence
null \cite{jo2026}. Recent work on behavioral entanglement measures how much the
behavior of one model can be predicted from another across six open-weight
families, but stops at an independence audit and verifier reweighting rather than
a trait-level variance partition \cite{kuai2026}.

\begin{table}[t]
\caption{\textbf{Differentiation against closest related work.}}
\label{tab:diff}
\setlength{\tabcolsep}{4pt}
\begin{tabular}{|p{92pt}|p{62pt}|p{110pt}|p{70pt}|}
\hline
Work & Object & What it does & Gap \\
\hline
Kim \emph{et al.}, ICML 2025 & agreement rate & measures when models co-err; cross-sectional & no variance partition; no temporal structure \\
Messing, TEE (2026) & pipeline facets & crossed random effects on judge/temperature/prompt; TEE-corrected CIs & facets are pipeline, not model population; no identifiability gate \\
Li \emph{et al.}, Tracing the Roots (ACL 2026) & dataset lineage graphs & reconstructs post-training data ancestry; contamination along paths & dataset lineage is a channel into the era component, not an error-trait partition \\
Li \emph{et al.}, Preference Leakage (ICLR 2026) & judge-model bias & judges favor related generators (same/inherit/family) & instrument-level bias, downstream of correlated error \\
PhyloLM / lineage reconstruction & phylogeny & reconstructs ancestry trees & no variance decomposition \\
Kuai \emph{et al.}, 2026 & behavioral entanglement & independence audit + verifier reweighting & entanglement indices, not a trait partition \\
Monoculture literature & deployed portfolio & welfare argument given correlated error & correlated error is an input, not the estimand \\
This work & error traits on connected subset & identifiability-gated $\sigma^{2}_{L}/\sigma^{2}_{E}/\sigma^{2}_{U}$ partition & --- \\
\hline
\end{tabular}
\end{table}

No verified prior work estimates $\sigma^{2}_{L}+\sigma^{2}_{E}+\sigma^{2}_{U}$ for
model error traits, on the provably connected subset, with the estimator's validity
established in simulation first. Every prior approach either documents correlation
without partitioning it, decomposes the variance of a different object, or treats
correlated error as a given input to a downstream argument.

\section{Formal Model and Identifiability}
\label{sec:model}

\subsection{Setup and notation}
Let the population consist of $N$ open-weight models indexed by $i\in\{1,\dots,N\}$.
Each model is released by a family $f(i)\in\{1,\dots,6\}$ (a verified design
grouping, e.g., Llama, Mistral, Qwen, Phi, Gemma, DeepSeek) in a release quarter
$e(i)\in\{1,\dots,14\}$ spanning 2023Q1--2026Q2. Each model receives a continuous
per-model trait $Y_{i}$ (Section \ref{sec:methods}), measured on a fixed common
item set. We model the trait as
\begin{equation}
Y_{i} = \mu + \alpha_{f(i)} + \beta_{e(i)} + u_{i},
\label{eq:model}
\end{equation}
where $\mu$ is the grand mean, $\alpha_{f}\sim\mathcal{N}(0,\sigma^{2}_{L})$ are
independent family (lineage) effects, $\beta_{e}\sim\mathcal{N}(0,\sigma^{2}_{E})$
are independent era effects, and $u_{i}\sim\mathcal{N}(0,\sigma^{2}_{U})$ are
independent model-specific residuals. Table~\ref{tab:notation} collects the
notation.

\begin{table}[t]
\caption{\textbf{Notation.}}
\label{tab:notation}
\setlength{\tabcolsep}{4pt}
\begin{tabular}{|p{52pt}|p{210pt}|}
\hline
Symbol & Meaning \\
\hline
$N$, $i$ & size of the open-model population (47); model index \\
$f(i)$, $e(i)$ & family and release quarter (era) of model $i$ \\
$F$, $E$ & number of families (6) and release quarters (14) \\
$Y_{i}$ & continuous per-model error trait, measured on the fixed common item set \\
$\mu$ & grand mean \\
$\alpha_{f}$, $\beta_{e}$, $u_{i}$ & family (lineage), era, and model-specific random effects in \eqref{eq:model} \\
$\sigma^{2}_{L}$, $\sigma^{2}_{E}$, $\sigma^{2}_{U}$ & lineage, era, and model-unique variance components \\
$\theta$ & variance-component vector $(\sigma^{2}_{L},\sigma^{2}_{E},\sigma^{2}_{U})$ \\
$\theta_{P}$ & observational share estimand (lineage conditional on era grouping) \\
$\theta_{M}$ & mechanistic share estimand (lineage with era held exactly fixed) \\
$\mathbf{C}$ & $N\times(F+E)$ design matrix of family and era indicators \\
$\mathbf{V}(\theta)$ & marginal covariance of $\mathbf{y}$ given \eqref{eq:cov} \\
$\ell_{R}(\theta)$ & log-restricted-likelihood (REML objective, \eqref{eq:reml}) \\
$\mathrm{VIF}$ & variance inflation factor of the design columns \\
\hline
\end{tabular}
\end{table}

In vector form,
\begin{equation}
\mathbf{y} = \mu\mathbf{1} + \mathbf{C}\boldsymbol{\gamma} + \mathbf{u}, \qquad
\boldsymbol{\gamma}\sim\mathcal{N}\bigl(\mathbf{0},\operatorname{diag}(\sigma^{2}_{L},
\sigma^{2}_{E})\bigr),
\label{eq:vector}
\end{equation}
The causal structure that motivates the two-estimand rule is summarized in
Fig.~\ref{fig:dag}: family membership influences the error trait directly and
through the release date, and release date is a mediator of the lineage path and
a potential confounder of the observational family--error association. Because
the mediational path cannot be separated from the lineage channel without holding
era fixed, a single causal attribution is impossible by construction.

\Figure[t!](topskip=0pt, botskip=0pt, midskip=0pt){figs/fig_dag.png}
{ \textbf{Causal structure of the error trait.} Family membership and era both
influence the error trait; release date is a mediator of the lineage path and a
potential confounder of the observational family--error association, motivating
the separate observational ($\theta_{P}$) and mechanistic ($\theta_{M}$)
estimands. Teacher leakage enters the era (shared environment) channel.
Appendix~\ref{app:identification} gives the formal identification
argument.\label{fig:dag}}

where $\mathbf{C}$ is the $N\times 12$ design matrix of family and era indicators
with $u_{i}$ as in \eqref{eq:model}. The marginal covariance is
\begin{equation}
\mathbf{V}(\theta) = \sigma^{2}_{U}\mathbf{I} + \mathbf{C}
\operatorname{diag}(\sigma^{2}_{L},\sigma^{2}_{E})\mathbf{C}^{\top},
\label{eq:cov}
\end{equation}
with $\theta=(\sigma^{2}_{L},\sigma^{2}_{E},\sigma^{2}_{U})$.

\subsection{Estimators}
We estimate $\theta$ by restricted maximum likelihood (REML), maximizing the
log-restricted-likelihood \cite{harville1977,patterson1971}
\begin{equation}
\ell_{R}(\theta) = -\tfrac{1}{2}\log\det\mathbf{V}
-\tfrac{1}{2}\log\det(\mathbf{1}^{\top}\mathbf{V}^{-1}\mathbf{1})
-\tfrac{1}{2}(\mathbf{y}-\mu\mathbf{1})^{\top}\mathbf{V}^{-1}
(\mathbf{y}-\mu\mathbf{1}),
\label{eq:reml}
\end{equation}
maximized directly with a numerical optimizer. Because $\mathbf{C}$ is low-rank
($F+E=20$ columns), the Woodbury identity
\begin{equation}
\mathbf{V}^{-1} = \sigma_{U}^{-2}\Bigl(\mathbf{I} -
\mathbf{C}\bigl(\sigma^{2}_{U}\mathbf{I} +
\operatorname{diag}(\sigma^{2}_{L},\sigma^{2}_{E})\mathbf{C}^{\top}\mathbf{C}\bigr)^{-1}
\mathbf{C}^{\top}\Bigr)
\label{eq:woodbury}
\end{equation}
reduces each likelihood evaluation to linear algebra on $20\times20$ matrices,
making repeated simulation feasible. Share confidence intervals use a Monte-Carlo
delta method on the log-variances with a numerically differentiated Hessian,
PSD-clipped for sparse designs.

We verified this direct implementation against the standard reference
implementations. On balanced data the REML estimates reproduce two-way ANOVA
method-of-moments estimates exactly (verified at 12$\times$12, 6$\times$14, and
8$\times$10 designs). In contrast, the widely used `statsmodels` crossed-variance
formulation does \emph{not} maximize the REML objective (Table~\ref{tab:estcomp}):
on a balanced 12$\times$12 design it returns a family/era split away from the
optimum (REML objective 66.674 vs.\ 65.994 for the direct optimizer), and on the
realistic D2 occupancy it understates the family share by roughly $5\times$ more
than the documented six-family small-sample limit. This is an upstream finding
worth reporting; all results here use the direct maximizer. The exact
comparison---optimizer settings, design, and the two objective values---is
reproduced in the released artifact
(src/lineage\_era/analysis/reml.py) and in the project decision log (2026-08-03);
a fresh reproduction of the family-share gap on D2 (scenario A) is included in
Table~\ref{tab:estcomp}.

\begin{table}[t]
\caption{\textbf{Estimator comparison: the direct maximizer vs.\ the
`statsmodels` crossed-variance path.} Balanced design: $F{=}E{=}12$, two models
per cell (288 models), fixed simulated dataset as in the decision log
(2026-08-03); the family-share bias column is measured on the realistic D2
occupancy (scenario A), reproduced 2026-08-06.}
\label{tab:estcomp}
\setlength{\tabcolsep}{4pt}
\begin{tabular}{|p{58pt}|p{30pt}|p{30pt}|p{30pt}|p{32pt}|p{42pt}|}
\hline
Method & Fam.\ $\sigma^{2}$ & Era\ $\sigma^{2}$ & Uniq.\ $\sigma^{2}$ & REML obj.\ $(-2\ell_{R})$ & Fam.-share bias on D2 (A) \\
\hline
Two-way ANOVA (method of moments) & 0.399 & 0.313 & 0.656 & 65.994 & --- \\
Direct REML (this work, Woodbury) & 0.399 & 0.313 & 0.656 & 65.994 & $-5.3$pp (300 reps)\\
`statsmodels` MixedLM ($vc\_formula$) & 0.609 & 0.476 & 0.656 & 66.674 & $-28.0$pp (40 reps)\\
\hline
\end{tabular}
\end{table}

\subsection{Estimands}
The primary observational estimand is the share vector
\begin{equation}
\theta_{P} = \left(\frac{\sigma^{2}_{L}}{\sigma^{2}_{L}+\sigma^{2}_{E}+\sigma^{2}_{U}},
\frac{\sigma^{2}_{E}}{\sigma^{2}_{L}+\sigma^{2}_{E}+\sigma^{2}_{U}},
\frac{\sigma^{2}_{U}}{\sigma^{2}_{L}+\sigma^{2}_{E}+\sigma^{2}_{U}}\right).
\label{eq:thetaP}
\end{equation}
The mechanistic estimand $\theta_{M}$ is the lineage share within designs that hold
era exactly fixed: co-released family cohorts and the verified cross-generation
fine-tune chains (the only documented parent--offspring edges), reported separately
from $\theta_{P}$.

\subsection{Identifiability conditions}
The partition \eqref{eq:thetaP} is estimable only if the design is crossed. We
require: (i) every family spans at least two release quarters and at least two
quarters contain at least two families; (ii) the induced design matrix has full
column rank; (iii) the variance-inflation factors of the design are bounded
($\mathrm{VIF}\le10$). A nested design---each family confined to a single era---is
perfectly aliased, and the identifiability gate must detect the aliasing rather
than silently ``succeed'' (Section \ref{sec:methods}).

\section{Methodology}
\label{sec:methods}
The protocol proceeds in three stages, each with an exit gate. The ordering is the
contribution: identifiability before results, simulation before real data, and
study-population design before empirical inference. Fig.~\ref{fig:pipeline}
shows the protocol spine; stages 1--3 are complete and reported here, and stage 4
is the pre-registered measurement pass whose outcomes are reported against the
pre-registered decision rule.

\Figure[t!](topskip=0pt, botskip=0pt, midskip=0pt){figs/fig_pipeline.png}
{ \textbf{Protocol architecture.} Three completed stages (population audit,
simulation validation, pre-analysis study-population design), each behind an
identifiability or recoverability gate, followed by the measurement and
decomposition stage (pre-registered eval pass, pending) and the decision layer.
The G3 stage is outcome-independent: it never observes trait values.
\label{fig:pipeline}}

\subsection{Study Population Construction}
\label{sec:phase0}
We assembled a verified population of $N=47$ open-weight general models spanning six
families and fourteen release quarters (2023Q1--2026Q2). The family $\times$ quarter
contingency table was built from Hugging Face API metadata and technical reports,
with documented divergences between the public release date and HF timestamps
corrected by hand. Two membership rules are central. First, \emph{all claims are
scoped to the connected subset}---the maximal subset on which lineage and era vary
jointly---so the partition \eqref{eq:thetaP} is defined. Second, hosted or
closed-weight models are excluded. The resulting design is crossed but unbalanced:
11 of 14 quarters contain at least two families and no family is confined to a
single era (Fig.~\ref{fig:design}).

\subsection{Simulation Validation of the Estimator}
\label{sec:phase1}
Before any real data are analyzed, the estimator is validated on simulated data
with known ground truth under three regimes \cite{harville1977}:
\begin{itemize}
\item \emph{D1, balanced crossed reference:} $30$ families $\times$ $14$ eras
  $\times$ $2$ models per cell; the large family count isolates estimator
  calibration from the small-sample limit of the real design.
\item \emph{D2, realistic occupancy:} The occupancy matrix copied from the study
  population table ($6$ families $\times$ $14$ quarters, $47$ models, sparse
  cells). Bias or collapse here is a gate failure: the real-data claim never
  proceeds.
\item \emph{D3, nested (must fail):} Each family confined to a single era, so
  $\sigma^{2}_{L}$ and $\sigma^{2}_{E}$ are perfectly aliased. Three independent
  detectors (BLUP family-vs-era collinearity, SE/estimate inflation, and
  profile-likelihood flatness) must flag the aliasing.
\end{itemize}
A liability decision was resolved in simulation before the trait path was fixed:
on item-level binary outcomes at the real occupancy, both a linear-probability
model and a binomial GLMM under-estimate the era component and drive it to the
boundary in 20--60\% of repetitions, whereas continuous per-model traits recover it
with 98--100\% coverage. Phase 2 therefore uses per-model continuous traits with
LPM--REML. A family $\times$ era interaction term was also simulated and documented
as non-identified at the sparse occupancy (SE/estimate ratios $\approx10^{4}$); the
interaction is excluded by design rather than absorbed.

\subsection{Pre-Analysis Study-Population Design}
\label{sec:g3}
The empirical campaign measures a fresh trait on all $47$ models. Because a full
run is expensive, we pre-registered a gate that selects the study population
\emph{before any GPU spend}: the smallest connected-subset population whose
family $\times$ quarter design is structurally identifiable \emph{and} whose era
recovery clears the strict Phase 1 bar in simulation. The key property of this
stage is that it is \emph{outcome-independent}: the optimizer's inputs are
occupancy (family $\times$ quarter), the lineage graph, identifiability
constraints, and cost---never trait values, accuracies, or evaluation outputs
(Table~\ref{tab:g3inputs}).

\begin{table}[t]
\caption{\textbf{G3 optimization inputs.}}
\label{tab:g3inputs}
\setlength{\tabcolsep}{4pt}
\begin{tabular}{|l|c|}
\hline
Input & Used \\
\hline
Family $\times$ quarter occupancy & $\checkmark$ \\
Lineage graph (verified edges, chain) & $\checkmark$ \\
Identifiability constraints (rank, VIF, span, crossing) & $\checkmark$ \\
Cost (public $>$ gated, est.\ GPU minutes) & $\checkmark$ \\
Trait values / accuracy & $\times$ \\
Error similarity & $\times$ \\
Evaluation outputs & $\times$ \\
\hline
\end{tabular}
\end{table}

The procedure is fully pre-registered. It minimizes model count subject to hard
constraints (all six families present with at least two quarters of span; every
quarter 2023Q1--2026Q2 retains at least one model; the in-subset verified-edge
endpoints and the documented chain are retained so the mechanistic estimand
survives; the induced design passes the structural checks), then tie-breaks on
cost. ``Valid,'' not merely ``identifiable,'' requires that the candidate design
recover known era shares in a fixed-design DGP battery (scenarios A and B) under
the strict bar (absolute era-share bias $\le5$pp, CI coverage $\ge90\%$,
convergence $\ge90\%$). Because the per-rep share bias is heavy-tailed at this
occupancy (standard deviation 22--26pp), a candidate that clears the bar at 300
repetitions must also clear it with a $\ge1$pp margin at a 1000-repetition
confirmation; knife-edge designs sitting exactly on the bar are rejected. Kept
models are assigned an auditable reason by single-model ablation---removing the
model and asking which design constraint fails---so the population file documents
\emph{why} each model is required (era-window, structural
identifiability, or statistical recoverability).

\subsection{Trait Measurement and Decomposition}
\label{sec:phase2}
On the selected population, the pipeline is: (i) a fresh five-shot MMLU pass
\cite{hendrycks2021} on the
common fixed item set (the only compliant source of per-question data; a public
reuse audit found no reusable per-question records); (ii) intake validation that
aborts on any contract violation of the manifest; (iii) assembly of the continuous
per-model trait; (iv) the identifiability gate on the measured design (hard fail =
abort); (v) the $\theta_{P}$ partition with REML \eqref{eq:reml}; (vi) bootstrap
confidence intervals combined with a trait-error Monte Carlo; and (vii) sensitivity
blocks (leave-one-family, leaked-drop, subject-drop, trait-definition, and a
documented sanity cross-check against prior leaderboards that is explicitly not a
validation). The era-convergence trend and the $\theta_{M}$ mechanistic tables are
reported as table entries, separately from the primary partition.

\section{Results}
\label{sec:results}
Results status: we report (a) the validated results that exist today---the
population audit, the simulation validation, and the
pre-analysis population design---and (b) the empirical outputs of the
measurement campaign (Section \ref{sec:results-empirical}), which are pending the
pre-registered eval pass and are explicitly marked as such. No empirical claim
about lineage-versus-era shares is made in this version; those claims are the
object of the pre-registered measurement pass. % TODO: replace the pending entries once the eval pass completes.

\subsection{Population Audit}
The verified population is $6$ families $\times$ $14$ quarters with $N=47$ models.
The design is crossed (unbalanced/incomplete): 11 of 14 quarters contain at least
two families, no family is confined to one era, and $5$ cross-generation
parent--offspring edges are verified from the metadata (the bulk of true lineage is
undocumented and drawn from technical reports). The connected subset---on which all
claims are scoped---is the population itself; a nested subpopulation exists and is
excluded by the scoping rule. Three caveats are carried into the analysis:
the Llama lineage terminates at Llama 4 (post-2025 era variation is carried by the
other families); DeepSeek V4 is a ground-up redesign treated as a new independent
lineage rather than a descendant; and cross-family teacher leakage (e.g., Phi-4
trained on GPT-4o-generated data) belongs in the era channel as shared environment.
Fig.~\ref{fig:design} shows the occupancy matrix.

\Figure[t!](topskip=0pt, botskip=0pt, midskip=0pt){figs/fig_design.png}
{ \textbf{Study population construction: family $\times$ release-quarter
occupancy of the 47-model connected subset.} 11 of 14 quarters contain at least
two families (highlighted) and no family is confined to a single era; sparse
cells carry the small-sample limits reported in Section \ref{sec:threats}. The
verified cross-generation lineage edges are shown in Fig.~\ref{fig:lineage}.
\label{fig:design}}

\Figure[t!](topskip=0pt, botskip=0pt, midskip=0pt){figs/fig_lineage.png}
{ \textbf{Verified lineage on the connected subset.} All 47 models on the
family $\times$ release-quarter grid (filled markers: gated access; open
markers: public). The five parent--offspring edges verified from the Hugging
Face \texttt{base\_model} field and the documented Mistral-Small chain support
the mechanistic estimand $\theta_{M}$; the primary partition $\theta_{P}$
uses the verified family grouping instead of the sparse edge record. Dashed
grey arrows mark the documented cross-family teacher-leakage channels, which
belong to the era (shared-environment) channel.\label{fig:lineage}}

\subsection{Simulation Validation}
\label{sec:results-phase1}
Table~\ref{tab:sim} summarizes the validation. Under the balanced reference (D1,
300 repetitions) share bias is $\le2.5$pp with coverage 95--96\%. Under the real
occupancy (D2) share bias is $\le5.3$pp with coverage 95--100\%; the family-share
bias of $-5.3$pp in the lineage-dominant scenario is the documented small-sample
limit of six family levels (df$=5$) and is reported as a limit, not corrected
away. The nested design (D3) is detected in 100\% of repetitions by all three
independent detectors, with zero silent coverage. The liability test resolved the
trait path to per-model continuous traits; the family $\times$ era interaction is
documented as non-identified rather than estimated. The verdict is GO WITH CHANGES,
with the changes being the continuous-trait path and the disclosure of the
family-share limit.

\begin{table}[t]
\caption{\textbf{Simulation validation (Phase 1).}}
\label{tab:sim}
\setlength{\tabcolsep}{4pt}
\begin{tabular}{|p{52pt}|p{95pt}|p{55pt}|p{45pt}|}
\hline
Regime & Recovery / detection & Gate & Verdict \\
\hline
D1 balanced-crossed & share bias $\le2.5$pp, coverage 95--96\% & PASS & GO \\
D2 realistic occupancy & share bias $\le5.3$pp, coverage 95--100\% & PASS$^{*}$ & GO (documented) \\
D3 nested (aliased) & detection 100\%, silent coverage 0\% & PASS & GO \\
Liability (binary) & era underpowered at real occupancy; continuous recovers & n/a & path decided \\
L$\times$E interaction & non-identified (SE ratio $\approx10^{4}$) & n/a & documented \\
\hline
\multicolumn{4}{p{240pt}}{$^{*}$The strict bar is defined on the \emph{era}-share
bias (mean over converged repetitions; Section \ref{sec:phase1}); the family-share
bias of $-5.3$pp in the lineage-dominant scenario is the six-family small-sample
limit (df$=5$) and is reported as a documented limit in Section
\ref{sec:threats}, not waived.}\\
\end{tabular}
\end{table}

\Figure[t!](topskip=0pt, botskip=0pt, midskip=0pt){figs/fig_d3.png}
{ \textbf{D3 aliasing detection under the nested (mis-specified) design.}
Percentage of repetitions in which each independent detector flags the
family--era aliasing, 300 repetitions per scenario on a fixed seed. A nested
design must \emph{fail detectably}---joint detection is 100\% and zero silent CI
coverage is observed---never silently ``succeed''. Detectors are defined formally
in Appendix~\ref{app:d3}.\label{fig:d3}}

\subsection{Pre-Analysis Population Design}
\label{sec:results-g3}
Fig.~\ref{fig:g3} and Table~\ref{tab:g3} report the gate. The structural minimum is
21 models (identifiable: full rank, VIF $\le10$), but the 21-model design sits
exactly on the strict bar in simulation: its era-share bias at 1000--2000
repetitions is $\approx-5.0$pp with a standard error of $0.6$--$0.8$pp---a
knife-edge whose verdict flips with the repetition count---so it fails the
margin confirmation. The minimum \emph{valid} population is 22 of 47 models,
which clears the bar at 300 repetitions (era bias A $2.4$pp / B $-0.8$pp), passes
the 1000-repetition margin confirmation (A $2.2$pp / B $-2.4$pp), and is robust at
2000 repetitions (A $1.7$pp / B $-3.2$pp). The full 47-model design passes by the
same criterion (A $1.5$pp / B $-2.2$pp at 1000 repetitions), confirming that the
reduction is statistically justified: 22 models recover era shares in the same
order of magnitude as all 47, at roughly one-third of the estimated single-GPU
cost (a 67\% reduction).
The winner keeps all six families, all fourteen quarters, the verified-edge
endpoints, and the chain needed for the mechanistic estimand. The population file
assigns each of the 22 kept models an auditable reason by single-model ablation
(era-window, structural identifiability, or statistical recoverability)
and each dropped model the reason ``redundant in-cell replication.''

\begin{table}[t]
\caption{\textbf{G3 gate: minimum valid population = 22 of 47.}}
\label{tab:g3}
\setlength{\tabcolsep}{4pt}
\begin{tabular}{|c|c|c|c|c|c|c|}
\hline
$n$ & models & bias A (pp) & bias B (pp) & cov\% (A/B) & confirm A/B & pass \\
\hline
47 & 47 & 0.3 & $-3.5$ & 99 / 96 & 1.5 / $-2.2$ & True \\
21 & 21 & 1.8 & $-4.9$ & 96 / 98 & 2.1 / $-5.1$ & False \\
22 & 22 & 2.4 & $-0.8$ & 98 / 99 & 2.2 / $-2.4$ & True \\
\hline
\end{tabular}
\end{table}

\Figure[t!](topskip=0pt, botskip=0pt, midskip=0pt){figs/fig_g3_trace.png}
{ \textbf{Pre-analysis population-design decision trace: population size vs.\
era-share bias (scenario B, confirmation battery) with the 2000-repetition
robustness points.} 47 passes, the structural minimum 21 fails the margin
confirmation (knife-edge), and 22 passes --- the extra model over the structural
minimum is statistically necessary, not computationally convenient. The winner
also clears the strict bar at 2000 repetitions (A $1.7$pp / B $-3.2$pp).
\label{fig:g3}}

\subsection{Empirical outputs (pending measurement)}
\label{sec:results-empirical}
The following results are produced by the Phase 2 pipeline on the measured trait
and are \emph{pending the fresh eval pass}; they are listed with the exact outputs
that will fill them. % TODO: replace all placeholders with the eval-pass outputs.
\begin{itemize}
\item \emph{Error-similarity panel (supporting).} Observed pairwise error overlap
  across models, chance-corrected with a pre-registered measure (phi is locked by
  the pre-registration rule), with the null ladder (observed $\to$
  matched-accuracy $\to$ item-difficulty $\to$ independence). Outputs: overlap
  heatmap, network, and embedding (Fig.~\ref{fig:similarity}).
\item \emph{{$\theta_{P}$} variance partition.} REML estimates and bootstrap
  intervals for $\sigma^{2}_{L}$, $\sigma^{2}_{E}$, $\sigma^{2}_{U}$ on the
  22-model population (Fig.~\ref{fig:partition}). \emph{(TBD)}.
\item \emph{Identifiability audit on the measured design.} Rank, VIF, and the
  D3-must-fail detectors re-run on the real occupancy. \emph{(TBD)}.
\item \emph{Era-convergence entry.} Era share by release quarter after lineage
  adjustment; reported as a table entry. \emph{(TBD)}.
\item \emph{{$\theta_{M}$} mechanistic tables.} Family contrasts within
  co-released quarters and per-quarter slopes along the verified fine-tune edges,
  reported separately. \emph{(TBD)}.
\item \emph{Sensitivity blocks.} Leave-one-family, leaked-drop, subject-drop,
  trait-definition, and the documented leaderboard sanity check. \emph{(TBD)}.
\end{itemize}

\Figure[t!](topskip=0pt, botskip=0pt, midskip=0pt){figs/fig_placeholder.png}
{ \textbf{Observed pairwise error overlap between models, chance-corrected
(pre-registered phi measure) with the null ladder.} \emph{Pending the fresh eval
pass.\label{fig:similarity}}}
% Placeholder; replace with the real overlap panel after the eval pass.

\Figure[t!](topskip=0pt, botskip=0pt, midskip=0pt){figs/fig_placeholder.png}
{ \textbf{{$\theta_{P}$} variance partition: lineage, era, and model-unique shares
with bootstrap intervals on the G3 minimum valid population (22 of 47), drawn
within the 47-model connected subset.} \emph{Pending the
fresh eval pass.\label{fig:partition}}}
% Placeholder; replace with the real share figure (variance_shares.pdf) after the eval pass.

\section{Discussion and Intervention Implications}
\label{sec:discussion}
The partition $\theta_{P}$ distinguishes two actionable regimes. If the lineage
share dominates, cross-family diversification is a credible mitigation: models with
independent ancestries add error variance that is not shared, so an ensemble is
less likely to fail a question the way its members do. If the era share dominates,
diversification alone is insufficient, because models released in the same period
share the same training environment regardless of ancestry. The decision rule that
follows from RQ6 is intentionally conservative: only a dominant era share refutes
diversification-as-remedy for the connected population; a dominant lineage share is
a necessary, not sufficient, condition for recommending diversification, because
the connected subset may not represent the broader model universe. Independent
recent work on ensemble combination reaches the same boundary from the other side:
pairwise error correlation substantially \emph{underestimates} the co-failure
ceiling of an ensemble, so the operative risk is the joint all-wrong rate
\cite{chen2026}; our partition estimates exactly the lineage and era shares whose
sum governs that ceiling on the connected population.

Two structural safeguards keep the discussion honest. First, the mechanistic
estimand $\theta_{M}$ is reported separately and never merged into $\theta_{P}$;
only the small co-released cohorts and staggered fine-tune chains can hold era
exactly fixed, so any mechanistic claim is scoped to them. Second, all conclusions
are scoped to the connected subset of the open-model population as defined in
Section \ref{sec:phase0}; the partition is undefined by construction on nested
subpopulations, which is a property of the design, not a data gap.

\section{Threats to Validity and Limitations}
\label{sec:threats}
\begin{itemize}
\item \emph{Six-family small-sample limit.} The real design has six family levels
  (df$=5$), which caps family-share CI coverage below nominal regardless of CI
  construction and produces a family-share bias of order $-5$pp in the
  lineage-dominant scenario (F2 in the simulation report). Family shares are
  reported with this limit disclosed; the same limit argues for acquiring more
  lineage families or reporting family shares as lower bounds.
\item \emph{Lineage metadata sparsity.} True cross-generation lineage is largely
  undocumented in model cards; only five cross-generation edges are verified.
  This constrains the mechanistic estimand, not the primary partition, which uses
  the verified family grouping.
\item \emph{Teacher leakage.} Cross-family teacher-student relationships (e.g.,
  Phi-4 from GPT-4o data, Gemma from Gemini) violate the independence assumption to
  a known-but-unknown degree and are disclosed as belonging to the era channel.
\item \emph{Underpowered binary era.} Item-level binary outcomes cannot resolve
  the era variance at this occupancy; the continuous per-model trait path is the
  documented resolution, and any binary-era claim must carry the underpower caveat.
\item \emph{Independent trait-error.} Per-model trait measurement errors are
  assumed independent across models; correlated measurement error would move shares
  jointly, and is probed by the trait-definition and leaked-drop sensitivity blocks.
\item \emph{Common item set.} A fresh five-shot pass on one fixed item set is the
  basis for all traits; prior work mixing leaderboard subsets is explicitly not used
  as validation.
\item \emph{Open-weight sampling frame.} The population is restricted to open-weight
  models with documented release records; hosted and closed models are excluded by
  design, so conclusions scope to the connected subset of the open-weight universe
  and do not generalize to closed systems.
\item \emph{Benchmark contamination and saturation.} Traits come from a widely cited
  five-shot benchmark; if pretraining included the items, trait variance is
  compressed toward zero and era/lineage shares are biased downward. The
  leaderboard sanity cross-check is explicitly not validation and is so labeled.
\item \emph{Power and boundary shares.} With $F=6$ families and 14 quarters, the
  era trend (RQ5) has limited support and share estimates near the boundary carry
  inflated CIs; these are reported with the underpower caveat rather than re-aimed.
\item \emph{Temporal drift.} The era component pools 14 quarters under an iid
  assumption; if the shared training environment is non-stationary within the
  window, the era effect is mis-specified. The era-convergence entry and the
  connected-window audit probe this assumption.
\item \emph{Visibility bias.} The frame is drawn from widely tracked open releases;
  smaller or niche families are underrepresented, which could overstate the family
  component relative to the full open-weight universe.
\item \emph{Off-subset scope.} The partition is defined only where lineage and era
  vary jointly; conclusions do not extend to nested subpopulations (where the
  partition is undefined by construction) or to closed models.
\end{itemize}

\section{Conclusion}
\label{sec:conclusion}
We have presented an identifiability-gated variance-decomposition instrument for
correlated errors in public language models, and the decision order that makes it
defensible: verify identifiability, validate the estimator in simulation, design the
study population without observing outcomes, and only then decompose measured
traits into lineage, era, and model-specific components on the connected subset.
The estimator is validated under both balanced and realistic occupancy and fails
detectably under nested mis-specification. A pre-analysis procedure selects a
minimum valid population of 22 of 47 models---the smallest population that is
structurally identifiable and clears the simulation-based recoverability bar---at
about one-third of the full-run cost (a 67\% reduction). The empirical partition---the actual lineage and
era shares on the selected population---is the outcome of the pre-registered
measurement pass and is deliberately not anticipated here; the decision rule for
whether cross-family diversification is a credible mitigation for correlated
failure is pre-registered and will be evaluated from those shares. All
artifacts---population, design, the optimizer, the validator, and the tests---are
released with the paper.

\appendices

\section{Phase 0 Verification}
Family groupings were verified against Hugging Face organization metadata and
technical reports. Release quarter is the public release date, not the HF
\emph{createdAt} timestamp; four documented divergences were corrected by hand.
DeepSeek V4 is treated as a new independent lineage (a ground-up redesign that
dropped the MLA architecture), and the within-family Small chain (Small
3 $\to$ 3.1 $\to$ 3.2 $\to$ 4; Devstral-2 on Small-3.1-Base) is the only verified
cross-generation chain.

\section{Formal Identification Argument}
\label{app:identification}
For the covariance model \eqref{eq:cov},
$\mathbf{V}(\theta)=\sigma^{2}_{U}\mathbf{I}+\sigma^{2}_{L}\mathbf{Z}_{F}\mathbf{Z}_{F}^{\top}
+\sigma^{2}_{E}\mathbf{Z}_{E}\mathbf{Z}_{E}^{\top}$. REML discards the fixed mean,
so $\theta$ is identifiable iff the three matrices
$\mathbf{I},\mathbf{Z}_{F}\mathbf{Z}_{F}^{\top},\mathbf{Z}_{E}\mathbf{Z}_{E}^{\top}$
are linearly independent on the orthogonal complement of $\mathbf{1}$. For
$i\neq j$, the $(i,j)$ entry of $\mathbf{V}$ is
$b\,\mathbf{1}\{f(i)=f(j)\}+c\,\mathbf{1}\{e(i)=e(j)\}$ with
$b=\sigma^{2}_{L}$, $c=\sigma^{2}_{E}$: a pair sharing only a family pins $b$, a
pair sharing only an era pins $c$, and a cross pair pins nothing. Failure modes
are then only (i) no family-sharing pair with disjoint eras exists---equivalently
family and era partitions coincide, the nested design---in which case
$\mathbf{Z}_{F}$ and $\mathbf{Z}_{E}$ span the same column space and
$\sigma^{2}_{L},\sigma^{2}_{E}$ are perfectly aliased; and (ii) the
family$\times$era incidence graph is disconnected, in which case the constraints
separate by component and components alias with a mixture. Every connected
crossed design with both marginals active is therefore identifiable. This is the
rank condition of Section~\ref{sec:model} and the gate of
Section~\ref{sec:methods}; Appendix~\ref{app:d3} probes it empirically.

\section{Detectable-Failure Battery (D3)}
\label{app:d3}
The nested mis-specification is flagged by three independent detectors, all
computed on 300 fresh repetitions of each scenario:
\begin{itemize}
\item \emph{BLUP collinearity.} The absolute correlation of family and era BLUPs
  ($|\operatorname{corr}(\hat{u}_{F},\hat{u}_{E})|$) exceeds $0.9$.
\item \emph{SE inflation.} Some variance component has
  $\widehat{\mathrm{SE}} \ge |\hat{\sigma}|$ (SE at least the estimate).
\item \emph{Profile flatness.} The drop in the profile likelihood over a
  $\pm0.4$ log-variance window along the aliased direction stays below
  $1.9207$ (the $\chi^{2}_{1}$ 95\% critical value over two).
\end{itemize}
Non-convergence is tracked as a fourth signal. A design passes the gate only if
all three detectors flag the aliasing in $\ge90\%$ of repetitions and no silent
CI coverage is observed; on the final design all three hit 100\% and silent
coverage is $0\%$.

\section{Crossing Implies Identifiable Rank}
Let the $N\times(F+E)$ incidence matrix be
$\mathbf{C}=[\mathbf{Z}_{F}\;\mathbf{Z}_{E}]$. The marginal covariance
\eqref{eq:cov} is $\mathbf{V}=\sigma^{2}_{U}\mathbf{I}+\mathbf{C}\mathbf{G}\mathbf{C}^{\top}$
with $\mathbf{G}=\operatorname{diag}(\sigma^{2}_{L}\mathbf{1}_{F},\sigma^{2}_{E}\mathbf{1}_{E})$.
The variance map $\theta\mapsto\mathbf{V}(\theta)$ is injective on the
orthogonal complement of $\mathbf{1}$ iff $\mathbf{C}$ has full column rank and
is not block-reducible. For a connected, crossed design with every family in
$\ge2$ eras and every era in $\ge2$ families, the contrast pairs in
Appendix~\ref{app:identification} exist for both channels, so the rank condition
holds; a nested or disconnected design reduces the effective rank to
$F$ (resp.\ separates components) and the map is non-injective. The Phase 1
design loop (Section~\ref{sec:g3}) searches exactly over these full-rank
crossed subpopulations; the winner (22 of 47) realizes rank 19 of 20.

\section{REML and Woodbury Computation}
The REML objective \eqref{eq:reml} is maximized over the three log-variances with
a numerical optimizer; each evaluation needs $\mathbf{V}^{-1}$ and
$\det\mathbf{V}$. Writing
$\mathbf{V}^{-1}=\sigma_{U}^{-2}(\mathbf{I}-\mathbf{C}(\sigma^{2}_{U}\mathbf{I}
+\mathbf{C}^{\top}\mathbf{C}\mathbf{G})^{-1}\mathbf{C}^{\top})$ (the Woodbury
identity, \eqref{eq:woodbury}), each likelihood evaluation reduces to linear
algebra on $(F+E)\times(F+E)=20\times20$ matrices---$\mathbf{C}^{\top}\mathbf{C}$
is built once---so a full evaluation is $O(N)$ with negligible constants.
Share confidence intervals use a Monte-Carlo scheme on the asymptotic
$3\times3$ covariance of the log-variances, with a PSD clip for sparse designs;
the optimizer, the validator, and every reported statistic are released, and the
estimator-comparison table (Table~\ref{tab:estcomp}) is reproducible from the
released code under the recorded seeds.

\section{Computation and Cost Plan}
All Phase 0 and Phase 1 analyses run on CPU and are fully reproducible. The
Phase 2 eval pass runs on one 8$\times$H200-141GB node in fp8; per-model
memory and GPU-time plans are released with the artifacts. Table~\ref{tab:cost}
is the cost plan for the selected 22-model population: 710 estimated
single-GPU minutes against 2{,}154 for the full 47-model roster---about one-third
of the full-run cost (a 67\% reduction). Eight of the 22 are gated and require
license acceptance.
\begin{table}[t]
\caption{Phase 2 GPU cost plan and access for the 22-model population
(cost column: estimated single-GPU minutes).}
\label{tab:cost}
\centering
\renewcommand{\arraystretch}{1.15}
\setlength{\tabcolsep}{4pt}
\scriptsize
\begin{tabular}{@{}llllll@{}}
\hline
Family & Model & Quarter & Params & Access & Est.\ min \\
\hline
Llama & Llama-1 & 2023Q1 & 7B & gated & 10 \\
Phi & Phi-1 & 2023Q2 & 1.3B & public & 2 \\
Qwen & Qwen-7B & 2023Q3 & 7B & public & 10 \\
Mistral & Mistral-7B & 2023Q3 & 7B & public & 10 \\
Phi & Phi-1.5 & 2023Q3 & 1.3B & public & 2 \\
Phi & Phi-2 & 2023Q4 & 2.7B & public & 4 \\
Qwen & Qwen1.5 & 2024Q1 & 72B & public & 102 \\
Phi & Phi-3 & 2024Q2 & 14B & public & 20 \\
Llama & Llama-3.1 & 2024Q3 & 70B & gated & 99 \\
Llama & Llama-3.3 & 2024Q4 & 70B & gated & 99 \\
Phi & Phi-4 & 2024Q4 & 3.8B & public & 5 \\
Mistral & Mistral-Small-3 & 2025Q1 & 24B & gated & 34 \\
Mistral & Mistral-Small-3.1 & 2025Q1 & 24B & public & 34 \\
Phi & Phi-4-reasoning-plus & 2025Q2 & 11B & public & 16 \\
Mistral & Mistral-Small-3.2 & 2025Q2 & 24B & gated & 34 \\
Gemma & Gemma-3n & 2025Q2 & 4B & gated & 6 \\
DeepSeek & DeepSeek-V3.1 & 2025Q3 & 671B/37B & public & 53 \\
Mistral & Devstral-2 & 2025Q4 & 24B & gated & 34 \\
DeepSeek & DeepSeek-V3.2 & 2025Q4 & 685B/37B & public & 53 \\
Phi & Phi-4-reasoning-vision-15B & 2026Q1 & 15B & public & 21 \\
Mistral & Mistral-Small-4 & 2026Q1 & 32B & gated & 45 \\
Gemma & Gemma-4-12B & 2026Q2 & 12B & public & 17 \\
\hline
\multicolumn{5}{@{}l}{22-model total} & 710 \\
\multicolumn{5}{@{}l}{47-model full roster} & 2{,}154 \\
\hline
\end{tabular}
\end{table}

\section{Minimum Valid Population Roster}
Table~\ref{tab:roster} records, for each of the 22 selected models, the single
binding reason for inclusion, matching the decision log. Fourteen are public;
eight are gated. Excluded models fail an era-window or a structural-identifiability
requirement, or sit off the connected subset.
\begin{table}[t]
\caption{Roster of the minimum valid population with the binding inclusion reason.}
\label{tab:roster}
\centering
\renewcommand{\arraystretch}{1.15}
\setlength{\tabcolsep}{2pt}
\scriptsize
\begin{tabular}{@{}llllll@{}}
\hline
Family & Model & Quarter & Params & Access & Reason \\
\hline
Llama & Llama-1 & 2023Q1 & 7B & gated & era-window \\
Phi & Phi-1 & 2023Q2 & 1.3B & public & era-window \\
Qwen & Qwen-7B & 2023Q3 & 7B & public & crossing \\
Mistral & Mistral-7B & 2023Q3 & 7B & public & rank/VIF \\
Phi & Phi-1.5 & 2023Q3 & 1.3B & public & rank/VIF \\
Phi & Phi-2 & 2023Q4 & 2.7B & public & era-window \\
Qwen & Qwen1.5 & 2024Q1 & 72B & public & era-window \\
Phi & Phi-3 & 2024Q2 & 14B & public & era-window \\
Llama & Llama-3.1 & 2024Q3 & 70B & gated & edge/chain \\
Llama & Llama-3.3 & 2024Q4 & 70B & gated & edge/chain \\
Phi & Phi-4 & 2024Q4 & 3.8B & public & edge/chain \\
Mistral & Mistral-Small-3 & 2025Q1 & 24B & gated & edge/chain \\
Mistral & Mistral-Small-3.1 & 2025Q1 & 24B & public & edge/chain \\
Phi & Phi-4-reasoning-plus & 2025Q2 & 11B & public & edge/chain \\
Mistral & Mistral-Small-3.2 & 2025Q2 & 24B & gated & edge/chain \\
Gemma & Gemma-3n & 2025Q2 & 4B & gated & crossing \\
DeepSeek & DeepSeek-V3.1 & 2025Q3 & 671B/37B & public & era-window \\
Mistral & Devstral-2 & 2025Q4 & 24B & gated & edge/chain \\
DeepSeek & DeepSeek-V3.2 & 2025Q4 & 685B/37B & public & edge/chain \\
Phi & Phi-4-reasoning-vision-15B & 2026Q1 & 15B & public & edge/chain \\
Mistral & Mistral-Small-4 & 2026Q1 & 32B & gated & edge/chain \\
Gemma & Gemma-4-12B & 2026Q2 & 12B & public & era-window \\
\hline
\end{tabular}
\end{table}
A model is \emph{era-window-forced} when removing it would strand an era window
with no usable model; \emph{edge/chain-forced} when it is the unique carrier of
a verified lineage edge or of the Mistral Small chain (the mechanistic
estimand's only cross-generation support); and \emph{crossing/rank-VIF-forced}
when it is required for a crossed full-rank design with $\mathrm{VIF}\le10$.

\section*{Acknowledgment}
% TODO: acknowledge funding, institutions, compute, and any contributors.

\begin{thebibliography}{10}

\bibitem{kim2025} E.~Kim, A.~Garg, K.~Peng, and N.~Garg, ``Correlated errors in
large language models,'' in \emph{Proc.\ Int.\ Conf.\ Mach.\ Learn.\ (ICML)},
vol.~267, pp.~30038--30066, 2025. [Online]. Available: \url{https://arxiv.org/abs/2506.07962}

\bibitem{harville1977} D.~A.~Harville, ``Maximum likelihood approaches to variance
component estimation and to related problems,'' \emph{J.\ Amer.\ Statist.\ Assoc.},
vol.~72, no.~358, pp.~320--338, 1977.

\bibitem{patterson1971} H.~D.~Patterson and R.~Thompson, ``Recovery of
inter-block information when block sizes are unequal,'' \emph{Biometrika},
vol.~58, no.~3, pp.~545--554, 1971.

\bibitem{brennan2001} R.~L.~Brennan, \emph{Generalizability Theory}. New York, NY,
USA: Springer, 2001.

\bibitem{phylolm2024} N.~Yax, P.-Y. Oudeyer, and S.~Palminteri, ``PhyloLM:
Inferring the phylogeny of large language models and predicting their performances
in benchmarks,'' in \emph{Proc.\ Int.\ Conf.\ Learn.\ Represent.\ (ICLR)}, 2025.
[Online]. Available: \url{https://arxiv.org/abs/2404.04671}

\bibitem{monoculture2024} B.~Hedden and M.~Raghavan, ``Algorithmic monoculture and
its critics,'' 2026. [Online]. Available: \url{https://arxiv.org/abs/2604.06047}

\bibitem{jo2026} N.~Jo, N.~Garg, and M.~Raghavan, ``The subjectivity of
monoculture,'' 2026. [Online]. Available: \url{https://arxiv.org/abs/2602.24086}

\bibitem{kuai2026} C.~Kuai \emph{et al.}, ``How independent are large language
models? A statistical framework for auditing behavioral entanglement and
reweighting verifier ensembles,'' 2026. [Online]. Available:
\url{https://arxiv.org/abs/2604.07650}

\bibitem{messing2026} S.~Messing, ``Hidden measurement error in LLM pipelines
distorts annotation, evaluation, and benchmarking,'' 2026. [Online]. Available:
\url{https://arxiv.org/abs/2604.11581}

\bibitem{li2026roots} Y.~Li \emph{et al.}, ``Tracing the roots: A multi-agent
framework for uncovering data lineage in post-training LLMs,'' in \emph{Proc.\
Annu.\ Meeting Assoc.\ Comput.\ Linguist.\ (ACL)}, 2026, pp.~9606--9625.
[Online]. Available: \url{https://arxiv.org/abs/2604.10480}

\bibitem{li2026pref} D.~Li \emph{et al.}, ``Preference leakage: A contamination
problem in LLM-as-a-judge,'' in \emph{Proc.\ Int.\ Conf.\ Learn.\ Represent.\
(ICLR)}, 2026. [Online]. Available: \url{https://arxiv.org/abs/2502.01534}

\bibitem{hendrycks2021} D.~Hendrycks \emph{et al.}, ``Measuring massive multitask
language understanding,'' in \emph{Proc.\ Int.\ Conf.\ Learn.\ Represent.\
(ICLR)}, 2021.

\bibitem{chen2026} J.~Chen, ``When does combining language models help? A
co-failure ceiling on routing, voting, and mixture-of-agents across 67 frontier
models,'' 2026. [Online]. Available: \url{https://arxiv.org/abs/2606.27288}

\end{thebibliography}

\begin{IEEEbiographynophoto}{S. Kanaga Suba Raja}% TODO: replace bio with real text
received the B.E. degree in Computer Science and Engineering and the M.E. degree
in Computer Science and Engineering, and is with the Department of Computer
Science and Engineering, SRM Institute of Science and Technology, Tiruchirappalli.
His research interests include ... . % TODO: fill interests
\end{IEEEbiographynophoto}

\begin{IEEEbiographynophoto}{Sudharsan S}% TODO: replace bio with real text
is with the Department of Computer Science and Engineering, SRM Institute of
Science and Technology, Tiruchirappalli. His research interests include ... . % TODO: fill interests
\end{IEEEbiographynophoto}

\EOD

\end{document}
